\chapter{Théorie du champ cristallin}
\label{chp:cristal_field}
Une première théorie pour décrire les propriétés des complexes, appelée théorie du champ cristallin, fut énoncée en 1929 par H. Bethe et H. Van Vleck. Il s'agit d'un modèle électrostatique dans lequel les ligands sont considérés comme des sphères dures chargées négativement dont l'agencement dans l'espace perturbe la structure électronique de l'ion métallique libre en considérant uniquement les orbitales $d$ de valence, les autres étant trop basses en énergie. Plusieurs conséquences découlent donc de ces hypothèses :
\begin{itemize}
\item La structure électronique du complexe dépend de sa géométrie ;
\item Il n’y a pas de mise en commun des électrons du métal et du ligand : on considère une liaison purement ionique.
\end{itemize}


\section{Construction de la théorie pour la géométrie octaédrique}

\begin{center}
\includegraphics[width=0.8\textwidth]{Fig01}
\captionof{figure}{Représentation d'un isocontour des $5$ orbitales $d$. Les couleurs indiquent les 2 phases d'une orbitale.}
\end{center}

Si on considère un atome, métal, n'ayant aucune interaction avec son environnement, les 5 orbitales $d$ sont dégénérés. Maintenant, si on plonge le métal dans un "bain" d'électron, puisque la densité électronique autour du métal est uniforme et isotrope, la répulsion des électrons du métal-électrons du bain conduit une déstabilisation uniforme. Les orbitales $d$ sont déstabilisées d'une même valeur d'énergie $\Delta E$.


\begin{center}
\includegraphics[width=0.8\textwidth]{Fig02}
\captionof{figure}{Variation d'énergie des orbitales $d$ d'un métal plongé dans un environnement électronique sphérique (champ sphérique ou uniforme), puis dans un environnement électronique octaédrique (champ octaédrique)}
\end{center}

Si on considère maintenant une géométrie octaédrique, la répulsion des orbitales $d$ du métal se fera de manière plus forte sur les axes $+x, -x, +y, -y, +z, -z$, où se trouve chacun des six ligands. Par conséquent, dans un champ octaédrique, les orbitales $d_{x^2 - y^2}$ et $dz^2$ sont plus destabilisées que dans un champ uniforme. Au contraire, les orbitales $d_{xy}$, $d_{yz}$ et $d_{xz}$ sont stabilisées par rapport au champ uniforme.


\definition{
La théorie du champ cristallin est un modèle électrostatique qui prédit que les orbitales $d$ d'un complexe métallique \textbf{ne sont pas} dégénérées. Le modèle de division des orbitales $d$ dépend du champ cristallin, qui est déterminé par l'\textbf{arrangement} et le \textbf{type} de ligand.
}

\section{Règle du barycentre et paramètre du champ cristallisation}

\definition{
La \textbf{règle du barycentre} stipule que la somme des énergies des orbitales $d$ dans un champ cristallin donné est conservée par rapport au champ cristallin uniforme. Ainsi,
\begin{align*}
\sum_i E(d_i) = 0
\end{align*}

~

On appelle le \textbf{paramètre de champ cristallin} octaédrique, la différence d'énergie $\Delta_O$ entre les niveaux $e_g$ ($dz^2$, $d_{x^2 - y^2}$) et $t_{2g}$ ($d_{xy}$,$d_{yz}$,$d_{xz}$).}

Par conséquent, avec les deux conventions au-dessus, on a deux équations à deux inconnues : 
\begin{align*}
\begin{cases} 3E(t_{2g}) + 2 E(e_g) &= 0 \\
E(e_g) - E(t_{2g}) &= \Delta_O 
\end{cases} \implies \begin{cases} E(t_{2g}) &= -\frac{2}{5} \Delta_O \\
E(e_g) &= \frac{3}{5} \Delta_O
\end{cases}
\end{align*}


On peut alors compléter le diagramme d'un élément du bloc d en géométrie (coordinence) octaédrique :

\begin{figure}[H]
\begin{minipage}{0.55\textwidth}
\includegraphics[width=\textwidth]{Fig03}
\end{minipage}\hfill\begin{minipage}{0.40\textwidth}
\captionof{figure}{\'{E}clatement des orbitales $d$ dans un champ cristallin octaédrique, où sont indiquées les différences d'énergie par rapport à l'état de référence sphérique.}
\end{minipage}
\end{figure}

\section{Autres géométries}

Une autre géométrie que l'on rencontrera souvent est la géométrie tétraédrique, où l'éclatement des orbitales est moindre (la distance ligand-métal est plus grande que dans un octaèdre équivalent) et inversé (ce sont les orbitales $t_2$ qui pointent vers les ligands).

\begin{center}
\includegraphics[width=0.8\textwidth]{Fig04}
\captionof{figure}{\'{E}clatement des orbitales $d$ dans la géométrie tétraédrique}
\end{center}

\`{A} cause de considérations purement géométriques, il existe une relation simple entre le paramètre de cristallin tétraédrique $\Delta_t$ et l'octaédrique $\Delta_O$ : 
\begin{align*}
\Delta_t = \frac{4}{9} \Delta_O
\end{align*}

Il est préférable de bien connaître le champ cristallin de ces deux géométries car tous les autres champs cristallins peuvent se déduire de celles-ci : 
\begin{center}
\includegraphics[width=0.8\textwidth]{Fig05}
\captionof{figure}{Quelques champs cristallins communs}
\end{center}

\section{Energie de stabilisation du champ cristallin (ESCC) et configuration électronique}
\definition{
L'\textbf{énergie de stabilisation du champ cristallin (ESCC)} est la différence d'énergie entre les électrons $d$ dans le champ cristallin réel par rapport aux électrons dans le champ sphérique.

~

La comparaison de deux ESCC permet de \textbf{comparer deux configurations électroniques} ou de \textbf{comparer deux géométries}, et d'identifier la plus stable (celle avec la plus petite ESCC).
}

Prenons par exemple l'atome \ce{Cr^{2+}} de configuration $d^4$. Celui-ci admet deux configurations possibles, représentées aux Fig. \ref{fig:06a} et \ref{fig:06b}, qui sont $(t_{2g})^4(e_g)^0$ et $(t_{2g})^3(e_g)^1$.

\begin{figure}[H]
\begin{minipage}{0.45\textwidth}
\subfloat[\label{fig:06a}]{\includegraphics[width=4cm]{Fig06a}}\subfloat[\label{fig:06b}]{\includegraphics[width=4cm]{Fig06b}}
\end{minipage}\hfill\begin{minipage}{0.50\textwidth}
\caption{Remplissage électronique pour un ion $d^4$ dans une \ref{fig:06a} configuration champ fort / bas spin ; \ref{fig:06b} configuration champ faible / haut spin.}
\end{minipage}
\end{figure}

L'énergie de stabilisation du champ cristallin permet de comparer ces deux configurations électroniques, qui dépendent de deux termes : 
\begin{enumerate}
\item le \textbf{coût énergétique} pour \textbf{apparier} deux électrons, appelé \textbf{énergie d'appariement} et notée $P$ ;
\item l'\textbf{éclatement} $\Delta_O$ entre les niveaux $e_g$ et $t_{2g}$.
\end{enumerate}

Calculons l'ESCC dans chacune des configurations : 
\begin{align*}
ESCC_a &= -4 \times \frac{2}{5} \Delta_O + P = -1.6 \Delta_O + P \\
ESCC_b &= -3 \times \frac{2}{5} \Delta_O + \frac{3}{5} \Delta_O = -0.6 \Delta_O 
\end{align*}

On a alors $\Delta ESCC = ESCC_b - ESCC_a = \Delta_O - P$, on en déduit alors deux situations possibles :
\begin{enumerate}
\item si $\Delta_O > P$, $ESCC_b > ESCC_a$, on parle alors de configuration "\textbf{champ fort}" ou "\textbf{bas spin}" ;
\item sinon, on parle de configuration "\textbf{champ faible}" ou "\textbf{haut spin}".
\end{enumerate}

\section{Paramètres influençant $\Delta_O$ et $P$}
On pourra retenir que le \textbf{paramètre de champ cristallin} $\Delta_O$ est influencé par : 
\begin{itemize}
\item la \textbf{nature du ligand} : il est difficile de rationaliser les tendances dans le cas des ligands, en général, plus l'atome lié au métal est "petit" et chargé (orbitales peu diffuses, forte densité électronique), plus $\Delta_O$ est grand. On a donc en général pour $\Delta_O$ : 

\begin{align*}
\ce{I^-}  < S^{2-}  < F^- < O^{2-}
\end{align*}

Il existe une série expérimentale, appelée \textbf{série spectrochimique} qui classe les ligands en fonction de la force du champ cristallin crée par sur un métal.

\item la \textbf{nature du métal} : plus le métal est diffus (vers la gauche et le bas dans le tableau périodique), plus $\Delta_O$ est grand ;
\item la \textbf{charge du métal} : plus le métal est chargé (degré d'oxydation élevé), plus $\Delta_O$ est grand.
\end{itemize}

Pour l'\textbf{énergie d'appariement}, on pourra retenir que l'énergie diminue lorsque les \textbf{orbitales} du métal sont plus \textbf{diffuses}. Par exemple, on aura : 
\begin{equation}
P( 3d ) > P( 4d ) > P( 5d )
\end{equation}


\definition{Au bilan, la théorie du champ cristallin permet par un approche intuitive pour décrire les comportements des complexes de métaux de transition des structures \textbf{ioniques}. Cependant, le modèle part du principe que l'on destabilise une structure en ajoutant un champ de ligand. Dans ce cas, d'où vient la stabilisation ? Bien entendu, que se passe-t-il si la liaison est \textbf{iono-covalente} avec une part non négligeable de \textbf{covalence} ? Il existe d'autres théories, plus complexes et abordées dans la suite de vos études, mais qui ne seront pas nécessaires au sein de ce module.}